\section*{Hoofdstuk \ref{ch:g1:living-or-not} --- Levend of niet-levend?}

\begin{solution}{exo:g1:living-or-not:1}
Geboren worden, eten, groeien, jongen kunnen krijgen, doodgaan.
Vandaag kun je een katje zien eten (melk drinken, happen) en zien
groeien --- het wordt week na week groter.
\end{solution}

\begin{solution}{exo:g1:living-or-not:2}
(a) Leeft --- hij is uit een ei gekropen, hij eet nectar, en er kunnen
babyrupsjes komen. (b) Leeft niet --- in een fabriek gebouwd, eet
nooit en groeit nooit. (c) Leeft --- het is uit een zaadje gegroeid.
(d) Leeft niet --- een wolk groeit en krimpt, maar ze is nooit geboren,
ze eet niet en ze krijgt geen jongen.
\end{solution}

\begin{solution}{exo:g1:living-or-not:3}
Nu leeft het niet meer. Het heeft wel geleefd: het groeide aan de boom
en werd door de boom gevoed. Het komt van een \omterm{def:g1:living-or-not:living}{levend wezen} --- de boom.
\end{solution}

\begin{solution}{exo:g1:living-or-not:4}
Stel de vragen van het leven samen: is de rivier geboren? Eet ze?
Groeit ze? Kan ze babyrivieren krijgen? Nee --- bewegen en lawaai
maken zijn op zichzelf geen kenmerken van leven. De rivier leeft niet,
ook al leven er heel veel levende wezens \emph{in}.
\end{solution}

\begin{solution}{exo:g1:living-or-not:5}
Leeft: de spin, de appel aan de boom. Heeft geleefd (of komt van een
\omterm{def:g1:living-or-not:living}{levend wezen}): de wollen sok (wol van een schaap), de gevallen
dennenappel (hij groeide aan een den). \omterm{def:g1:living-or-not:nonliving}{Nooit geleefd}: de regendruppel.
\end{solution}

\begin{solution}{exo:g1:living-or-not:6}
Nee. De vlam zakt voor de vragen van het leven: hij is aangestoken,
niet geboren, en vooral: er zijn geen babyvlammen die opgroeien --- een
vlam kan wel naar een andere kaars overgezet worden, maar hij krijgt
geen jongen zoals een slak of een madeliefje. (Het is een eerlijke
vraag: de vlam doet twee kenmerken echt goed na!)
\end{solution}

\begin{solution}{exo:g1:living-or-not:7}
Bijvoorbeeld: een houten stoel (van een boom), een katoenen T-shirt
(van een katoenplant), een wollen deken (van een schaap), een leren
schoen (van een koe), honing in de kast (van bijen).
\end{solution}

\begin{solution}{exo:g1:living-or-not:8}
De boon verandert: hij zwelt op, dan komt er een klein wit worteltje
uit, en daarna een groen scheutje. De kiezel blijft precies hetzelfde.
Dat laat zien dat de boon leefde --- een stil, wachtend \omterm{def:g1:living-or-not:living}{levend wezen} ---
en dat de kiezel niet leeft.
\end{solution}

\begin{solution}{exo:g1:living-or-not:9}
Is het geboren? De slak is uit een ei gekropen; de robot is gebouwd.
Kan het jongen krijgen? Er kunnen babyslakjes komen; babyrobots
bestaan niet. (Eten werkt ook: van sla groeit de slak, van de batterij
groeit de robot niet.)
\end{solution}

\begin{solution}{exo:g1:living-or-not:10}
Een paddenstoel wordt geboren (hij verschijnt uit piepkleine begintjes
in de grond), hij eet (hij haalt zijn voedsel uit de grond en uit dode
bladeren), hij groeit --- soms in één enkele nacht --- hij maakt nieuwe
paddenstoelen, en hij gaat dood. Alle kenmerken van leven zijn er: het
is een \omterm{def:g1:living-or-not:living}{levend wezen}, ook zonder poten en zonder groene bladeren.
\end{solution}
